%auto-ignore
%      this ensures the arxiv doesn't try to start TeXing here.
%!TEX root = super_lattice_models_draft.tex
%      prev line helps TeXShop do the right thing

%%%%%%%%%%%%%%%%
\section{Introduction}
%%%%%%%%%%%%%%%%

The renormalization group provides a method for understanding how a field theory arises as the continuum limit of a many-body lattice model. It gives both quantitative and qualitative information on which microscopic degrees of freedom and couplings remain relevant, and so enter into the effective long-distance description. In its half-century of existence, it has transformed the way theorists think about both statistical mechanics and field theory.

It is important to remember, however, that the renormalization group is not the only method for understanding how continuum arises from lattice. In this paper we go deeper than effective field theory, as we exploit the remarkable confluence between integrable lattice models, conformal field theory, and mathematics. In particular, we utilise the profound connection between knot and link invariants \cite{Jones1987} and lattice statistical mechanics \cite{Jones1989,Wadati1989,Wu1992}, a line of development that had started well before \cite{Temperley1971,Fortuin1971,Baxter1976}. 
Since this classic work, substantial progress has been made in developing the mathematics involved. In particular, much more is known about {\em fusion categories}. Among other things, fusion categories give a precise way of manipulating graphs while preserving their topological invariants  \cite{MSReview89,barrett1999,kup_spider,blob_paper,jones_pa,Kitaev2006}. Our aim is to show how the topological framework not only provides an elegant setup for lattice statistical mechanics, but gives a perspective both qualitatively and quantitatively useful. 

We use the inherent topological structure of a wide class of two-dimensional lattice models to define {\em topological defect lines} with remarkable properties. The partition function of a lattice model defined using a fusion category is independent of deformations of the paths of any topological defects present. Moreover, the construction results in a host of exact linear identities relating partition functions with different defect configurations. The reason the topological approach is so powerful is that many constraints are necessary to enforce consistency. The solutions of these constraints are simple to express in terms of topological data, whereas a brute-force approach typically is horrendously complicated.  

Universal results can be obtained simply and exactly from lattice computations.
One main result is our demonstration how lattice topological defects implement {\em dualities}, extending the work of Kramers and Wannier \cite{Kramers1941} well beyond the Ising and Potts models where it is familiar.
We hope to convince that these defects provide the best way of understanding why and how dualities work for two-dimensional classical lattice models and for quantum spin chains. The key observation is that distinct models can be defined using the same fusion category. The dualities then arise in a very natural fashion, as the topological defects built using the category separate regions described by the different models. When the model is the same on both sides, the self-duality places strong constraints on it, allowing us to compute universal quantities.  



The topological defects we construct inherit properties from the fusion category used to define them, providing a set of linear identities relating partition functions involving different defect configurations.  In essence, the microscopic definitions of the lattice model and its Boltzmann weights result in macroscopic constraints on the partition functions.  
A basic example is
\begin{align}
 \BubbleMacroLeft \; =\; \delta_{RB} \sqrt{\frac{d_P d_G}{d_R}} \; \BubbleMacroRight 
\label{bubble}
\end{align}
where $R,G,P,B$ label the different types of defects. The numbers $d_a$ are given by the category, and are known as {quantum dimensions}. Roughly speaking, the weight of a single defect line with label $a$ is $d_a$. 
A more general set of relations are known as $F$-moves.  They describe a change of basis, relating partition functions with distinct defect configurations as
\begin{align}
\FmoveMacroLeft \ =\  \sum_{Y} \big[ F_{P}^{RGB} \big]_{XY}\FmoveMacroRight\quad .
\label{MacroF}
\end{align}
The coefficients are called $F$-symbols, and are a glorified version of the $6j$ symbols familiar from quantum mechanics.
Finding a consistent set of $F$-symbols satisfying \eqref{MacroF} results in a huge number of non-linear relations.
The category hands us the solution.

Relations like \eqref{bubble} and \eqref{MacroF} are exact on the lattice. When a lattice model we study has a continuum limit, such partition function identities become exact relations for the corresponding field-theory quantities. These relations allow us to compute exact universal quantities in terms of topological invariants. When the lattice model is critical and ensuing field theory is conformally invariant, these quantities can be compared to those computed using  conformal field theory (CFT) \cite{BPZ:CFT:84}. One type of universal quantity we derive exactly and rigorously on the lattice is the ratio of $g$-factors for conformal boundary conditions \cite{Cardy1989,AffleckLudwig}. Another is the shift in momentum arising from twisting the boundary conditions using a topological defect. This quantity is related to the eigenvalues of a Dehn twist, and hence the conformal spin in the corresponding CFT \cite{Cardy1986}. The relations \eqref{bubble} and \eqref{MacroF} themselves correspond nicely to those arising directly in CFT \cite{Petkova2001,Frohlich:2006,Frohlich2009}. 

The connection of our approach to CFT brings the story full circle. Fusion categories were prefigured in the statistical mechanics of models having CFTs describing their continuum limit \cite{Temperley1971,Fortuin1971,Baxter1976}.  
Canonical examples of such lattice models are the Ising and Potts models  \cite{Baxter1982,Wu1992} and their ``parafermion'' generalizations \cite{Fradkin1980,Fateev82}, along with the integrable height models of Andrews, Baxter and Forrester \cite{Andrews1984}.  Characterizing the topological properties of fields in rational CFT provided a crucial impetus for the introduction of modular tensor categories (fusion categories with more data, including braiding)  \cite{MSReview89}. Further connections of CFT to topology and knot invariants came through the Chern-Simons approach to the Jones polynomial \cite{Witten:1989}. Our work exploits the subsequent progress in the understanding of fusion categories to derive new properties of the very same lattice models.  
We emphasize, however, that even though integrable models and conformal field theory are deeply embedded in the story, our method yields many results for non-integrable and non-critical models as well, duality being one prominent example. 


Despite this paper being a sequel, it is self-contained. In \cite{Aasen2016} (henceforth referred to as Part~I), we defined the Ising model in traditional fashion, and explained how topological defects had some rather magical properties. Underlying all these results were fusion categories. The magic remains in the much more general results described here, but now we display our tricks.

We are well aware that many readers may be unfamiliar with the topological approach to lattice models. We therefore have written three review sections describing the setup. In section \ref{sec:Fusion_Theory} we give a concise introduction to fusion categories, focusing on how to compute topological invariants. In section \ref{sec:statmech}, we explain how to build lattice height models using these categories. The Boltzmann weights of Potts, parafermion and ABF height models, among others, then can be related to topological invariants. In section \ref{modules} we explain the mathematics underlying our construction of lattice topological defects, the Turaev-Viro-Barrett-Westbury state sum.

Sections \ref{sec:topological_defect_lines} and \ref{sec:trivalent} contain our core results. The definition of lattice topological defect lines and the proof they satisfy the defect commutation relations is given in section \ref{sec:topological_defect_lines}. In section \ref{sec:trivalent}, we define topological junctions of defects, and prove that defect lines themselves satisfy the rules of the fusion category. In our opinion, the brevity of these two sections demonstrates the substantial value of the categorical approach.

We then give applications. In section \ref{sec:duality} we implement Kramers-Wannier duality using topological defects, and generalize it to a much larger class of models. Section \ref{sec:selfduality} explores a particular consequence of self-duality, enabling the derivation of degeneracies in the ground and low-lying states of certain non-critical models. We show in section \ref{sec:bdry} how our setup provides a natural way of describing boundary conditions in terms of states. In the critical case, it gives a way of computing universal ratios of $g$-factors easily on the lattice. In section \ref{sec:twists}, we define twisted boundary conditions in terms using topological defects and derive identities for the resulting partition functions. By understanding the behavior under Dehn twists, we explain how to extract exact results for conformal scaling dimensions as well. The conclusions are given in section \ref{sec:conclusion}. 






